\documentclass[10pt,a4paper]{amsart}
\usepackage[margin=19mm]{geometry}
\usepackage{amsmath,amssymb,mathtools}
\usepackage[scr=rsfs,cal=euler]{mathalfa}
\usepackage{xfrac}
\usepackage[unicode,hidelinks,bookmarks=false]{hyperref}
\newcommand{\ie}{i.e.\ }
\newcommand{\cf}{cf.\ }
\newcommand{\resp}{respectively\ }
\newcommand{\Subsection}{\subsection}
\providecommand{\nobreakdash}{\nobreak\hbox{-}}
\setlength{\emergencystretch}{2em}
\multlinegap0pt
\usepackage{bm}
\usepackage{bbm}
\usepackage{dsfont}
\usepackage{xspace}
\usepackage{cancel}
\usepackage{mathtools}
\usepackage{amsthm}
\makeatletter
\@ifundefined{thm}{\newtheorem{thm}{Thm}}{}
\@ifundefined{cor}{\newtheorem{cor}[thm]{Corollary}}{}
\@ifundefined{prop}{\newtheorem{prop}[thm]{Proposition}}{}
\makeatother
\title[Quantitative Thomas-Yau uniqueness]{Quantitative Thomas-Yau uniqueness}
\author{Yang Li}
\date{Source: arXiv:2207.08047v1 (CC BY 4.0)}
\begin{document}
\maketitle

\noindent\emph{Source and licence.} Quantitative Thomas-Yau uniqueness. Version arXiv:2207.08047v1 (CC BY 4.0), distributed under the Creative Commons Attribution 4.0 International licence, as recorded in the arXiv licence metadata for this exact version and shown on the abstract page https://arxiv.org/abs/2207.08047v1. Full rights notice: licenses/arxiv-2207.08047-CC-BY-4.0.txt.\par\smallskip

\noindent\textit{Editorial scope.} Counted unit: the Prop whose source label is \texttt{monotonicity} in the article (source0.tex lines 642--656, source label \texttt{monotonicity}), delivered together with the article's own complete original proof of it (source0.tex lines 782--818, one \texttt{proof} environment of 37 lines). No mathematics is written, repaired or re-derived: statement and proof are the article's own text line for line. The proof is detached from the statement in the source: the article states the result at lines 642--656 and prints its proof at lines 782--818, and the proof environment's own header names the statement, so the association is the article's own. Required context retained uncounted: HalfHausdorff (lines 571--576). Every definition, display and label of those lines is retained unchanged. Omitted from this package and not redistributed: 1--570, 577--641, 657--781, 819--1138. Nothing inside a retained excerpt is omitted or altered: every retained body is delivered exactly as the article's lines are written, so no omission receipt of any kind accompanies this package. Citations: the retained excerpts carry no citation command at all, so no bibliography entry of the article is transcribed and no reference is redistributed. Reference adaptation: none was needed and none was made; every reference inside the delivered unit resolves inside it, so no unresolved reference or citation is produced. Printed slips kept literal: no printed word, symbol, hypothesis or number of the counted statement or of its proof was altered. Layout and class: the article's own class is \texttt{article} with packages including babel, inputenc, amsmath, graphicx, amssymb, amsthm, tikz-cd, mathrsfs, todonotes, enumitem, yfonts, dsfont, MnSymbol, slashed; the wrapper is the lane's pinned \texttt{amsart} editorial wrapper at 10pt on A4, which restates the theorem-like environments the retained text needs and adds the article's own macro definitions that the retained text needs (none), with extra wrapper packages (bm, bbm, dsfont, xspace, cancel, mathtools). The delivered numbering is the unit's own. Assets: no retained span contains a figure environment, a graphics-inclusion command, a PDF-page inclusion, a table or a TikZ picture; no image file is copied into this package, the package declares no asset dependency and no image file is redistributed with it. Licence: version arXiv:2207.08047v1 of this article is distributed under the Creative Commons Attribution 4.0 International licence, as recorded in the arXiv licence metadata for this exact version and shown on the abstract page https://arxiv.org/abs/2207.08047v1. A full rights notice is delivered at licenses/arxiv-2207.08047-CC-BY-4.0.txt. Characters: every retained excerpt is pure ASCII, so the delivered engine, XeLaTeX with its default Latin Modern text font, reports no missing character.
\begin{cor}\label{HalfHausdorff}
Under the same conditions,
\begin{equation}
\sup_{p\in L} \text{dist}(p,L') \leq C\epsilon^{1/2n}.
\end{equation}
\end{cor}

\begin{prop}\label{monotonicity}
	Let $L'$ be a smooth Lagrangian with $|\theta|\leq \frac{\pi}{2}-\epsilon_0$, and 
	there exists $p'\in L'$ with $|p'|\leq \lambda\ll 1$. We shall use $C(\epsilon_0)$ to denote constants depending only on $\epsilon_0$ and the coordinate chart.
\begin{enumerate}
	\item 
For the Euclidean balls $B_{Eucl}(r)$  with radius $r>\lambda$ contained in the chart, we have the volume lower bound $\text{Vol}(L'\cap B_{Eucl}(r))\geq C(\epsilon_0)(r-\lambda)^n$.
	
	\item Assume furthermore that $\text{Vol}( \{ \theta_{L'}> \epsilon   \}\cap L' ) \leq \lambda^n $ where $\lambda\geq \epsilon^2$. Then for  the Euclidean balls $B_{Eucl}(r)$  with radius $r>C(\epsilon_0)\lambda$ contained in the chart, we have the sharper volume lower bound
\begin{equation}
\text{Vol}(L'\cap B_{Eucl}(r)) \geq \omega_n (r-C(\epsilon_0)\lambda)^n (1-C(\epsilon_0)r).
\end{equation}
\end{enumerate}


\end{prop}

\begin{proof}
We can take $\lambda=C\epsilon^{1/2n}$, so that by Cor. \ref{HalfHausdorff}, any $p\in L$ lies within distance $\lambda$ to $L'$. Our setting implies $\lambda\geq \epsilon^2$ and $\text{Vol}(\{\theta_{L'}>\epsilon  \})\leq \lambda^n$. By Prop. \ref{monotonicity}, and the fact that Euclidean balls and Riemannian balls only differ by relative error $O(r)$, we see that for $O(1)\geq r\geq C\lambda$, 
\[
\text{Vol}_g (L'\cap B_g(p,r)) \geq \omega_n (r-C\lambda)^n (1-Cr),\quad \forall p\in L.
\]
The constants are uniform for $p\in L$, so after integration
\[
\begin{split}
& \int_L e^{-\rho(p)}\text{Vol}_g (L'\cap B_g(p,r)) dvol_L(p) 
\\
\geq  & \omega_n (r-C\lambda)^n (1-Cr)\int_L e^{-\rho}dvol_L 
\\
\geq & \omega_n (r-C\lambda)^n (1-Cr)\int_L \text{Re}\Omega.
\end{split}
\]
The LHS can be computed by Fubini's theorem as
\[
\int_{L'} e^{-\rho(q)}dvol_{L'}(q) \int_{L\cap B_g(q,r)} e^{\rho(q)-\rho(p)} dvol_L(p).
\]
The function $\rho$ is smooth, so $|\rho(q)-\rho(p)|\leq Cr$. The Lagrangian $L$ has fixed $C^1$-regularity bound, so for small radius $r\ll 1$,
\[
\int_{L\cap B_g(q,r)} dvol_L(p)\leq (1+Cr) \omega_n r^n.
\]
Combining the above,
\[
\begin{split}
& \omega_n (r-C\lambda)^n (1-Cr)\int_L \text{Re}\Omega\leq \int_L e^{-\rho(p)}\text{Vol}_g (L'\cap B_g(p,r)) dvol_L(p) 
\\
\leq & (1+Cr) \omega_n r^n \int_{L'\cap \{   \text{dist}(\cdot, L)\leq r \} } e^{-\rho(q)}dvol_{L'}(q),
\end{split}
\]
whence for $C\lambda\leq r\leq O(1)$,
\begin{equation}\label{distanceweakL1}
\int_{L'\cap \{   \text{dist}(\cdot, L)\leq r \} } e^{-\rho(q)}dvol_{L'}(q) \geq (1- Cr- C\frac{\lambda}{r}) \int_L \text{Re}\Omega. 
\end{equation}
The choice $r= \epsilon^{1/4n}$ yields the claim.
\end{proof}

\end{document}
